\item \begin{enumerate}
    \item Determine la densidad de masa de un protón, representándolo como una esfera sólida de radio $1.00 \times 10^{-15}\text{ m}$.
    \item \textbf{¿Qué pasaría si?} Considere un modelo clásico de un electrón como una esfera sólida con la misma densidad que el protón. Determine su radio clásico.
    \item Imagine que este electrón posee un momento angular orbital de espín $L_{\text{espín}} = \hbar/2$ debido a su rotación clásica alrededor del eje $z$. Determine la rapidez lineal de un punto en el ecuador del electrón y
    \item compare cuantitativamente esta rapidez con la rapidez de la luz $c$.
\end{enumerate}